\scanpage{126}
Si $T$ est de rang impair, donc $T\simeq V\oplus\mathcal O_S\ldots$, on définira \deleted{\unsure{une}} $\sigma$ par
\[
\sigma(x)=x-2\pi(x)a\tag*{(22)}
\]
ou encore, sur $E$, par la symétrie p.r. à l'origine $\unsure{a|0}$. On aura encore (21).

3. \uline{Schéma des autom d'une quadrique affine} \inserted{birégulière}.

Soit
\[
\begin{aligned}
G&=\underline{\operatorname{Aut}}(P,H,Q)=\underline{\operatorname{Aut}}\,\mathrm{aff}(E,Q\cap E)\\
&\deleted{\simeq\operatorname{Aut}(T,\unsure{q},\pi\unsure{)},a,q)}
\end{aligned}
\]
\[
\deleted{(\pi\in\Gamma(S,\check T),a\in\Gamma(S,T),q(a)=1,\varphi(a,x)=}
\]
Dans le cas où $T$ de rang impair i.e. $E$ de rang pair, on a
\[
G\simeq\operatorname{Aut}(V,q_0)\qquad V\ \text{de rang pair}
\]
et désignant par $G^+$ le sous-schéma en groupes qui opère trivialement sur $\det V$, on trouve
\[
\deleted{\unsure{G\simeq G^+\times\mu_2|G\simeq G^+\times\mathbb Z/2\mathbb Z}}
\]
\diagram{\[
\begin{tikzcd}[ampersand replacement=\&,column sep=large]
\unsure{0|\cdot|1} \arrow[r,phantom,"{\unsure{-|\longrightarrow}}" description] \& G^+ \arrow[r] \& G \arrow[r] \& \mathbb Z/2\mathbb Z \arrow[r,phantom,"{\unsure{-|\longrightarrow}}" description] \& 0
\end{tikzcd}
\]}
Si $\unsure{T|E}$ de rang pair, \inserted{$G\simeq\operatorname{Aut}(\unsure{T|E},q,a)$} on va montrer qu'on a un iso can.
\[
G\simeq G^0\times(\mathbb Z/2\mathbb Z)_S
\]
où $G^0\xrightarrow{\sim}\underline{\operatorname{Aut}}^+(V,q_0)$, $(\mathbb Z/2\mathbb Z)_S$ engendré par $\sigma$.
% archive sheet number: 125
